\paperchapter{December 2019}{29 December 2019}{A}{50}{14th GSET}
\keytable{\kcell{1}{A} & \kcell{2}{D} & \kcell{3}{A} & \kcell{4}{C} & \kcell{5}{C} & \kcell{6}{A} & \kcell{7}{B} & \kcell{8}{C} & \kcell{9}{A} & \kcell{10}{D} \\
\kcell{11}{D} & \kcell{12}{C} & \kcell{13}{B} & \kcell{14}{C} & \kcell{15}{D} & \kcell{16}{A} & \kcell{17}{D} & \kcell{18}{B} & \kcell{19}{A} & \kcell{20}{D} \\
\kcell{21}{Z} & \kcell{22}{A} & \kcell{23}{D} & \kcell{24}{D} & \kcell{25}{D} & \kcell{26}{B} & \kcell{27}{D} & \kcell{28}{A} & \kcell{29}{B} & \kcell{30}{D} \\
\kcell{31}{B} & \kcell{32}{B} & \kcell{33}{C} & \kcell{34}{B} & \kcell{35}{C} & \kcell{36}{B} & \kcell{37}{A} & \kcell{38}{D} & \kcell{39}{A} & \kcell{40}{D} \\
\kcell{41}{C} & \kcell{42}{D} & \kcell{43}{C} & \kcell{44}{C} & \kcell{45}{C} & \kcell{46}{C} & \kcell{47}{B} & \kcell{48}{A} & \kcell{49}{B} & \kcell{50}{A}}
\begin{question}{1}
{Which set of the following differentiate the offline with that of online testing method ? The statements are:\newline (1) Offline is local centered, online is centralized\newline (2) Offline is non-digital, online is digital\newline (3) Offline is subjective, online is objective\newline (4) Internet is must for online, and not a necessity for offline\newline (5) Online usually provides instant evaluation, offline need not\par}
\optsfour{1, 4 and 5}{1, 2 and 5}{1, 3 and 5}{2, 3 and 4}
\answer{A}{1, 4 and 5}
\begin{explain}
The features that really distinguish online testing from offline testing are:
\begin{itemize}
\item (1) offline testing is \emph{local} (conducted at the centre), while online testing is \emph{centralised} (server-based);
\item (4) the internet is a must for online testing but not for offline testing;
\item (5) online testing usually gives \emph{instant} evaluation, while offline testing need not.
\end{itemize}
Statement (2) does not hold, since offline tests may also be digital (e.g.\ OMR). Statement (3) is false, since both modes can be subjective or objective. The answer is 1, 4 and 5.
\end{explain}
\end{question}

\begin{question}{2}
{The strength of project method lies in:\par}
\optstwo{Social participation}{Collection of information}{Role of teacher as instructor}{Individualised engagement}
\answer{D}{Individualised engagement}
\begin{explain}
In Kilpatrick's \textbf{project method} (based on Dewey's pragmatism) each learner \emph{plans, executes and evaluates} a purposeful activity in a real-life setting. Its strength lies in the learner's \textbf{individualised, active engagement}: learning by doing. The teacher acts as a guide rather than an instructor.
\end{explain}
\end{question}

\begin{question}{3}
{Computer based testing is essential to:\par}
\optsone{Make testing transparent, objective and efficient}{Ease teacher from burden of evaluation}{Prepare digitised evaluation sheets for speedy results}{Minimise human bias in testing}
\answer{A}{Make testing transparent, objective and efficient}
\begin{explain}
Computer-based testing (CBT) makes testing \textbf{transparent, objective and efficient}: standardised administration, randomised question sets, automatic scoring, quick results and secure records. The other options are narrower by-products of this main purpose.
\end{explain}
\end{question}

\begin{question}{4}
{Open book examination is to:\par}
\optsone{do away from malpractice of copying}{test ability of candidates to refer books and learn}{to test candidates ability on application based knowledge}{to increase the quality of question paper}
\answer{C}{to test candidates ability on application based knowledge}
\begin{explain}
An \textbf{open-book examination} lets candidates consult books. It therefore does not test memory; it tests the \textbf{ability to apply knowledge}: analysis, interpretation and problem solving.
\end{explain}
\end{question}

\begin{question}{5}
{By reflective teaching learner :\par}
\optsone{updates skills to be competent}{will be a good achiever in his examination}{will internalise knowledge and be judgemental}{will get content mastery}
\answer{C}{will internalise knowledge and be judgemental}
\begin{explain}
\textbf{Reflective teaching} makes learners think about what they learn and how. They \textbf{internalise} knowledge and develop the ability to \emph{judge} and evaluate critically (the reflective level of learning in Hunt's model: memory $\rightarrow$ understanding $\rightarrow$ reflective).
\end{explain}
\end{question}

\begin{question}{6}
{The data derived through experimental method:\par}
\optstwo{is not from a natural setting}{lacks specific measurement}{is unrealistic}{is artificial}
\answer{A}{is not from a natural setting}
\begin{explain}
In the \textbf{experimental method} the researcher manipulates independent variables and controls extraneous ones, often in a laboratory. The data therefore come from a \textbf{contrived, not a natural setting}. This is its main limitation (low ecological validity).
\why{Experimental data actually involve \emph{precise} measurement, so (B) is wrong. ``Unrealistic'' and ``artificial'' overstate the point; the defining feature is the controlled rather than natural setting.}
\end{explain}
\end{question}

\begin{question}{7}
{Which one of the following is a constraint for adopting historical method of research ?\par}
\optstwo{Lack of technical advancement}{Non-availability of primary data}{Ambiguity of concepts}{Subjective information}
\answer{B}{Non-availability of primary data}
\begin{explain}
Historical research depends on records of the past. Its biggest constraint is the \textbf{non-availability of primary sources} (original documents, eyewitness accounts), which are often lost, destroyed or inaccessible. Researchers then have to rely on secondary sources.
\end{explain}
\end{question}

\begin{question}{8}
{After completion of work, on a research problem, the report was discussed in a scholar's meet, and researcher was suggested to change title of problem and hypothesis, based on the data analysis. This is:\par}
\optstwo{Logical and appropriate}{Processes systematisation of report}{Illogical and not desirable}{Improving the validity of report}
\answer{C}{Illogical and not desirable}
\begin{explain}
The title and hypotheses are formulated \emph{before} data collection and guide the entire study. Changing them after the data have been analysed, to fit the results, amounts to HARKing (Hypothesising After the Results are Known). It is \textbf{illogical and not desirable} and undermines the scientific integrity of the research.
\end{explain}
\end{question}

\begin{question}{9}
{Title of the thesis is :\par}
\optsfour{An assertion}{An observation}{A fact}{An elaboration}
\answer{A}{An assertion}
\begin{explain}
A research title is a precise statement of what the study is about: the variables, population and relationship being investigated. As per the official key, the title is treated as an \textbf{assertion}, a declarative statement of the research problem.
\end{explain}
\end{question}

\begin{question}{10}
{Given below are the steps of a scientific research. Arrange them in order:\newline (1) Formulation of hypothesis\newline (2) Development of research tool\newline (3) Identification of research problem\newline (4) Selection of sample\newline (5) Data analysis and interpretation\par}
\optsfour{1, 2, 3, 5, 4}{1, 4, 3, 2, 5}{3, 4, 1, 2, 5}{3, 1, 2, 4, 5}
\answer{D}{3, 1, 2, 4, 5}
\begin{explain}
The logical sequence of scientific research:
\begin{center}
\begin{tikzpicture}[node distance=3mm,every node/.style={draw=brand,rounded corners,fill=soft,font=\scriptsize,align=center,text width=15mm,minimum height=9mm}]
\node (a) {(3) Identify problem}; \node[right=of a] (b) {(1) Formulate hypothesis}; \node[right=of b] (c) {(2) Develop tool}; \node[right=of c] (d) {(4) Select sample}; \node[right=of d] (e) {(5) Analyse \& interpret};
\foreach \x/\y in {a/b,b/c,c/d,d/e} \draw[-Stealth,brand] (\x)--(\y);
\end{tikzpicture}
\end{center}
The order is \textbf{3, 1, 2, 4, 5}.
\end{explain}
\end{question}

\begin{context}
\textbf{Directions for Q.~11 to 15:} Read the following passage and answer the questions.\par\smallskip
If the world is different today, we need to think, talk and engage accordingly. Falling back is unlikely to help, he said adding ``purposeful pursuit of national interest is shifting global dynamics''. Highlighting India's new approach in dealing with terrorism he compared ``lack of response to the terror attack of Mumbai in comparison to the way the country responded to the Uri and Pulwama strikes. On India walking away from the Regional Comprehensive Economic Partnership, the external affair minister said no agreement was better than having a bad agreement.

Giving a historical perspective to geopolitical issues, Jai-Shankar said ``For years India's position on the world state seemed assured but the 1962 conflict with China significantly damaged India's standing.
\end{context}
\begin{question}{11}
{RCEP stands for:\par}
\optsone{Regional Common Economic Party}{Regional Common Election Party}{Regional Comprehensive Election Partnership}{Regional Comprehensive Economic Partnership}
\answer{D}{Regional Comprehensive Economic Partnership}
\begin{explain}
\textbf{RCEP} is the \textbf{Regional Comprehensive Economic Partnership}, a free-trade agreement among the 10 ASEAN countries and China, Japan, South Korea, Australia and New Zealand (signed 2020). India withdrew from the negotiations in November 2019.
\end{explain}
\end{question}

\begin{question}{12}
{What was country's idea about Economic Partnership?\par}
\optsone{All Agreements are good}{Agreements are superior than country}{Nothing wrong in saying no to unfavourable agreement than having an awful agreement}{Say yes to all agreement}
\answer{C}{Nothing wrong in saying no to unfavourable agreement than having an awful agreement}
\begin{explain}
The minister said that ``\textbf{no agreement was better than having a bad agreement}''. Option (C) expresses this: there is nothing wrong in saying no to an unfavourable agreement rather than accepting an awful one.
\end{explain}
\end{question}

\begin{question}{13}
{The above statement was given by:\par}
\optstwo{Subhramanyam Swami}{Subhramanyam Jayshankar}{Ravi Shankar}{K. Subhramanyam}
\answer{B}{Subhramanyam Jayshankar}
\begin{explain}
The passage names \textbf{(S.) Jaishankar}, India's External Affairs Minister, as the speaker. Option (B) (``Subhramanyam Jayshankar'') refers to Dr.~Subrahmanyam Jaishankar.
\end{explain}
\end{question}

\begin{question}{14}
{India's New Approach on Terrorism is:\par}
\optsone{Wait till any action takes place}{Deny for any type of Action}{Earliest it was reactive, at present it is proactive}{Ready to loose for Peace}
\answer{C}{Earliest it was reactive, at present it is proactive}
\begin{explain}
The passage contrasts the ``lack of response'' to the 2008 Mumbai attack with India's response to the Uri (2016, surgical strikes) and Pulwama (2019, Balakot air strike) attacks. The approach has changed from \textbf{reactive to proactive}.
\end{explain}
\end{question}

\begin{question}{15}
{Year's back conflict of India with China was:\par}
\optstwo{National Issue}{International Issue}{Political Issue}{Geopolitical Issue}
\answer{D}{Geopolitical Issue}
\begin{explain}
The minister was ``giving a historical perspective to \textbf{geopolitical} issues'' when he spoke of the 1962 conflict with China, so the conflict is described as a geopolitical issue.
\end{explain}
\end{question}

\begin{question}{16}
{During the process of communication :\par}
\optsone{Sender encodes for the media and receiver decodes the message}{Sender decodes for the media and receiver encodes the message}{Medium encodes and decodes the message from sender}{Medium encodes and decodes the message from receiver}
\answer{A}{Sender encodes for the media and receiver decodes the message}
\begin{explain}
In the communication process the \textbf{sender encodes} the message into symbols suited to the medium, and the \textbf{receiver decodes} it to extract the meaning. The medium only carries the message.
\end{explain}
\end{question}

\begin{question}{17}
{Proxemics is the type of non-verbal communication which deals with:\par}
\optsfour{Appearance}{Touch}{Body organs}{Space}
\answer{D}{Space}
\begin{explain}
\textbf{Proxemics} deals with the use of \textbf{space} and distance in communication.
\why{Appearance belongs to artifactual communication; touch is haptics; movements of body organs are kinesics.}
\end{explain}
\end{question}

\begin{question}{18}
{Language barriers in communication exist due to :\par}
\optsone{the difference in the sender's recognition vocabulary and the receiver's user vocabulary}{the difference in sender's user vocabulary and the receiver's recognition vocabulary}{the sender's user vocabulary and the receiver's recognition vocabulary being alike}{the sender's recognition vocabulary and the receiver's user vocabulary being the same}
\answer{B}{the difference in sender's user vocabulary and the receiver's recognition vocabulary}
\begin{explain}
Every person has a larger \emph{recognition} vocabulary (words understood) than \emph{use} vocabulary (words actually used). A language barrier arises when the sender's \textbf{use vocabulary} contains words that are not in the receiver's \textbf{recognition vocabulary}, i.e.\ the two differ.
\end{explain}
\end{question}

\begin{question}{19}
{In intra-personal communication :\par}
\optsone{sender sends and receives the message}{sender sends and another receiver receives the message}{sender sends and two receivers receive the message}{sender sends and the group receives the message}
\answer{A}{sender sends and receives the message}
\begin{explain}
In \textbf{intrapersonal communication} a person communicates \emph{with oneself} (thinking, self-talk, reflection), so the same person is both sender and receiver.
\end{explain}
\end{question}

\begin{question}{20}
{The mass media impacts our society in different ways except:\par}
\optstwo{for influencing public opinion}{for modification of behaviour}{for changing values and belief system}{for judicial decisions}
\answer{D}{for judicial decisions}
\begin{explain}
Mass media shape public opinion, modify behaviour and change values and beliefs. \textbf{Judicial decisions} must be based on law and evidence and should not be influenced by the media, so this is the exception.
\end{explain}
\end{question}

\begin{question}{21}
{Find out the missing number: 4, 9, 20, 43, \ldots{}\par}
\optsfour{84}{96}{108}{133}
\dropped{}
\begin{explain}
The pattern is $\times2+1, \times2+2, \times2+3, \ldots$:
\[4\xrightarrow{\times2+1}9\xrightarrow{\times2+2}20\xrightarrow{\times2+3}43\xrightarrow{\times2+4}\mathbf{90}\]
90 is not among the options, so the agency \textbf{dropped} the question.
\end{explain}
\end{question}

\begin{question}{22}
{From Vadodara station, if we buy two tickets to station P and three tickets to station Q, the total cost is Rs. 77, but if we buy 3 tickets to station P and 5 tickets to station Q, the total cost is Rs. 124. The fare from Vadodara to P is:\par}
\optsfour{Rs. 13}{Rs. 14}{Rs. 18}{Rs. 20}
\answer{A}{Rs. 13}
\begin{explain}
$2P+3Q=77$ and $3P+5Q=124$.
Multiply the first equation by 5 and the second by 3: $10P+15Q=385$ and $9P+15Q=372$. Subtracting gives $P=\mathbf{13}$ (and $Q=17$).
Check: $2(13)+3(17)=77$ \checkmark; $3(13)+5(17)=124$ \checkmark.
\end{explain}
\end{question}

\begin{question}{23}
{300 gm of sugar solution has 40\% sugar in it. How much sugar should be added to make it 50\% in the solution ?\par}
\optsfour{30 gm}{40 gm}{50 gm}{60 gm}
\answer{D}{60 gm}
\begin{explain}
Sugar $=40\%\times300=120$ g. Let $x$ g of sugar be added.
\[\frac{120+x}{300+x}=\frac12\Rightarrow240+2x=300+x\Rightarrow x=\mathbf{60\ g}\]
\end{explain}
\end{question}

\begin{question}{24}
{Numbers in boxes given below follow the same arrangement. Find the missing numbers in the third box.\par}
\qfigure{\pyr{29}{16}{13}{27}{9}{33}\qquad\pyr{38}{21}{17}{35}{16}{43}\qquad\pyr{55}{30}{25}{?}{?}{?}}
\optsfour{65, 36, 63}{53, 24, 63}{49, 32, 53}{51, 25, 61}
\answer{D}{51, 25, 61}
\begin{explain}
Two rules hold in every box:
\begin{itemize}
\item \textbf{Bottom middle} $=$ (difference of the middle row)$^2$: $(16-13)^2=9$, $(21-17)^2=16$, so $(30-25)^2=\mathbf{25}$.
\item \textbf{Bottom left $+$ bottom right} $=2\times\text{top}+2$: $27+33=60=2(29)+2$, $35+43=78=2(38)+2$, so the sum must be $2(55)+2=112$. Also bottom right $-$ bottom left $=6, 8, \ldots$, so $\mathbf{10}$.
\end{itemize}
Left $+$ right $=112$ and right $-$ left $=10$ give \textbf{51 and 61}. The missing numbers are \textbf{51, 25, 61}.
\end{explain}
\end{question}

\begin{question}{25}
{By selling a mobile for Rs. 1,536, Amit lost 20\%. What percent shall he gain or loose by selling it for Rs. 2,000?\par}
\optstwo{A loss of 8.25\%}{A profit of 8.25\%}{A loss of 4.17\%}{A profit of 4.17\%}
\answer{D}{A profit of 4.17\%}
\begin{explain}
CP $=\dfrac{1536}{0.80}=1920$. New SP $=2000$, so profit $=80$.
\[\text{Profit \%}=\frac{80}{1920}\times100=\mathbf{4.17\%}\]
\end{explain}
\end{question}

\begin{question}{26}
{The Sanskrit word `Upamana' means:\par}
\optsfour{Implication}{Comparison}{Inference}{Perception}
\answer{B}{Comparison}
\begin{explain}
\textbf{Upamana} is knowledge by \textbf{comparison} or analogy. For example, a person told that ``a gavaya (wild ox) resembles a cow'' recognises a gavaya on seeing one.
\tip{Pratyaksha = perception; Anumana = inference; Upamana = comparison; Shabda = verbal testimony; Arthapatti = postulation; Anupalabdhi = non-perception.}
\end{explain}
\end{question}

\begin{question}{27}
{Odometer is to milage as compass is to:\par}
\optsfour{Speed}{Hiking}{Needle}{Direction}
\answer{D}{Direction}
\begin{explain}
An odometer is the instrument that measures \emph{mileage} (distance travelled); a compass is the instrument that shows \textbf{direction}. The relation is \emph{instrument : what it measures}.
\end{explain}
\end{question}

\begin{question}{28}
{Which of the following diagrams correctly represents Lions, Elephants and Animals?\par}
\optsfour{\tikz[baseline=-0.5ex,scale=0.3]{\draw (0,0) circle(1.3); \draw (-0.55,0) circle(0.5); \draw (0.55,0) circle(0.5);}}{\tikz[baseline=-0.5ex,scale=0.3]{\draw (0,0) circle(0.8); \draw (1.1,0) circle(0.8); \draw (2.2,0) circle(0.8);}}{\tikz[baseline=-0.5ex,scale=0.3]{\draw (0,0) circle(1.1); \draw (0,0.1) circle(0.5); \draw (1.6,0) circle(0.9);}}{\tikz[baseline=-0.5ex,scale=0.3]{\draw (0,0) circle(0.8); \draw (1.1,0) circle(0.8); \draw (3,0) circle(0.8);}}
\answer{A}{\tikz[baseline=-0.5ex,scale=0.3]{\draw (0,0) circle(1.3); \draw (-0.55,0) circle(0.5); \draw (0.55,0) circle(0.5);}}
\begin{explain}
Lions and elephants are both \textbf{animals}, and no lion is an elephant. So the two small circles must be \textbf{separate} and both lie \textbf{inside} the large ``Animals'' circle, as in diagram (A).
\begin{center}
\begin{tikzpicture}[scale=0.7,font=\scriptsize]
\draw[fill=brand!8] (0,0) circle(2.1); \node at (0,1.6) {Animals};
\draw[fill=white] (-0.95,-0.2) circle(0.85) node {Lions}; \draw[fill=white] (0.95,-0.2) circle(0.85) node {Elephants};
\end{tikzpicture}
\end{center}
\end{explain}
\end{question}

\begin{question}{29}
{Nyaya Syllogism has \blank{} statement.\par}
\optsfour{One}{Five}{Three}{Two}
\answer{B}{Five}
\begin{explain}
The Nyaya syllogism (\emph{Panchavayava}) has \textbf{five} members: Pratijna (proposition), Hetu (reason), Udaharana (example), Upanaya (application) and Nigamana (conclusion).
\end{explain}
\end{question}

\begin{question}{30}
{Should high chimneys be installed in industries?\newline Arguments:\newline (I) Yes, it reduces pollution at ground level.\newline (II) No, it increases pollution in upper atmosphere.\par}
\optstwo{Only argument I is strong}{Only argument II is strong}{Either I or II is strong}{Both I and II are strong}
\answer{D}{Both I and II are strong}
\begin{explain}
Both arguments are \textbf{strong}. Tall chimneys disperse pollutants high up and so reduce ground-level concentration, which directly protects people (I). But they spread the pollution into the upper atmosphere and over larger areas, contributing to acid rain and similar problems (II). Each argument is important and directly related to the question.
\end{explain}
\end{question}

\begin{question}{31}
{The sum of squares of three numbers is 138, while sum of their products taken two at a time is 131. Their sum is:\par}
\optsfour{16}{20}{24}{26}
\answer{B}{20}
\begin{explain}
\[(a+b+c)^2=(a^2+b^2+c^2)+2(ab+bc+ca)=138+2(131)=400\ \Rightarrow\ a+b+c=\mathbf{20}\]
\end{explain}
\end{question}

\begin{question}{32}
{Given the following wage distributions:\par}
\qfigure{\begin{tabular}{lc}\toprule \textbf{Wages (Rs.)} & \textbf{No.\ of workers}\\\midrule
Less than 50 & 4\\ Less than 60 & 10\\ Less than 70 & 20\\ Less than 80 & 40\\ Less than 90 & 50\\\bottomrule\end{tabular}}
What percentage of workers earn wage more than Rs.~70?
\optsfour{50}{60}{65}{75}
\answer{B}{60}
\begin{explain}
This is a cumulative (``less than'') distribution. Total workers $=50$, and those earning less than Rs.~70 number 20.
Workers earning more than Rs.~70 $=50-20=30$, and $\dfrac{30}{50}\times100=\mathbf{60\%}$.
\end{explain}
\end{question}

\begin{question}{33}
{Study of the following graph carefully and answer the question:\par}
\qfigure{\begin{tikzpicture}
\begin{axis}[ybar,width=11cm,height=4.6cm,bar width=12pt,ymin=0,ymax=28000,
  symbolic x coords={Gujarat,Haryana,Karnataka,Maharashtra,Rajasthan,Tamil Nadu,U.P.,Others},xtick=data,
  x tick label style={rotate=35,anchor=east,font=\scriptsize},yticklabel style={font=\scriptsize},
  nodes near coords,every node near coord/.append style={font=\tiny},
  title={\scriptsize Production of Rose in Various States},axis x line*=bottom,axis y line*=left,
  scaled y ticks=false]
\addplot[fill=white,draw=ink] coordinates {(Gujarat,15000) (Haryana,12500) (Karnataka,12450) (Maharashtra,20000) (Rajasthan,12400) (Tamil Nadu,22500) (U.P.,22400) (Others,25000)};
\end{axis}
\end{tikzpicture}}
Which of the following state(s) contribute(s) less than 10 per cent in the total Rose production?
\optsone{Rajasthan and Karnataka}{Only Rajasthan}{Rajasthan, Karnataka, Harayana}{Rajasthan, Karnataka, Haryana and Gujarat}
\answer{C}{Rajasthan, Karnataka, Harayana}
\begin{explain}
Total $=15000+12500+12450+20000+12400+22500+22400+25000=142{,}250$, so 10\% $=14{,}225$.
States below 14,225: \textbf{Haryana} (12,500), \textbf{Karnataka} (12,450) and \textbf{Rajasthan} (12,400). Gujarat (15,000) is above 10\%.
\end{explain}
\end{question}

\begin{question}{34}
{Study the pie-chart carefully to answer the question:\par}
\qfigure{\begin{minipage}{0.45\linewidth}\centering
{\scriptsize\bfseries Percentage of students (Total 4500)}\\[2pt]
\begin{tikzpicture}\pie[radius=1.35,text=pin,font=\tiny,color={brand!10,brand!30,brand!15,brand!40,brand!20,brand!5}]{11/Finance,22/HR,16/Mktg,18/IB,12/Banking,21/IT}\end{tikzpicture}
\end{minipage}\hfill
\begin{minipage}{0.45\linewidth}\centering
{\scriptsize\bfseries Ratio of Men to Women}\\[2pt]\small
\begin{tabular}{lcc}\toprule Specialization & Men & Women\\\midrule
HR & 2 & 3\\ Marketing & 7 & 5\\ IB & 5 & 1\\ Banking & 17 & 10\\ IT & 1 & 2\\ Finance & 5 & 4\\\bottomrule\end{tabular}
\end{minipage}}
The number of women studying Finance is what percentage of the number of men studying the same?
\optsfour{70}{80}{85}{90}
\answer{B}{80}
\begin{explain}
Only the Finance ratio is needed: Men : Women $=5:4$.
\[\frac{\text{Women}}{\text{Men}}\times100=\frac45\times100=\mathbf{80\%}\]
(Numerically: Finance $=11\%$ of 4500 $=495$, of whom men $=275$ and women $=220$; $220/275=80\%$.)
\end{explain}
\end{question}

\begin{question}{35}
{The distribution of weights (in kg) of 50 students in a class is as follows:\par}
\qfigure{\begin{tabular}{cc}\toprule \textbf{Weight (kg)} & \textbf{No.\ of students}\\\midrule
35--38 & 3\\ 38--41 & 13\\ 41--44 & 19\\ 44--47 & 10\\ 47--50 & 5\\\bottomrule\end{tabular}}
What is the average weight of the class?
\optsfour{40.26}{41.88}{42.56}{43.25}
\answer{C}{42.56}
\begin{explain}
Use class mid-points $x$:
\begin{center}\small
\begin{tabular}{cccc}\toprule Class & $x$ & $f$ & $fx$\\\midrule
35--38 & 36.5 & 3 & 109.5\\ 38--41 & 39.5 & 13 & 513.5\\ 41--44 & 42.5 & 19 & 807.5\\ 44--47 & 45.5 & 10 & 455.0\\ 47--50 & 48.5 & 5 & 242.5\\\midrule
& & 50 & 2128.0\\\bottomrule\end{tabular}
\end{center}
Mean $=2128/50=\mathbf{42.56}$ kg.
\end{explain}
\end{question}

\begin{question}{36}
{`Internet Explorer' is a type of:\par}
\optstwo{Operating System}{Browser}{IP Address}{Compiler}
\answer{B}{Browser}
\begin{explain}
Internet Explorer (Microsoft, 1995--2022) is a \textbf{web browser}, software for accessing and viewing websites. Its successor is Microsoft Edge.
\end{explain}
\end{question}

\begin{question}{37}
{Which of the following is an open source software?\par}
\optstwo{Mozilla Firefox}{Acrobat Reader}{MS PowerPoint}{Windows}
\answer{A}{Mozilla Firefox}
\begin{explain}
\textbf{Mozilla Firefox} is free and open-source (released under the Mozilla Public License), so its source code is publicly available to study and modify.
\why{Acrobat Reader is freeware but proprietary; MS PowerPoint and Windows are proprietary commercial software.}
\end{explain}
\end{question}

\begin{question}{38}
{NMEICT stands for:\par}
\optsone{National Mission on Environment through ICT}{National Mission on English through ICT}{National Mission on Economics through ICT}{National Mission on Education through ICT}
\answer{D}{National Mission on Education through ICT}
\begin{explain}
\textbf{NMEICT} is the \textbf{National Mission on Education through Information and Communication Technology}, launched by MHRD in 2009. It supports e-content (e.g.\ NPTEL, e-PG Pathshala), connectivity and the SAKSHAT portal.
\end{explain}
\end{question}

\begin{question}{39}
{Dr. Ramesh is Assistant Professor in one college and wishes to do refresher course in his subject. Which of the following will be useful to him?\par}
\optsfour{ARPIT}{MOODLE}{NPTEL}{ePG Pathshala}
\answer{A}{ARPIT}
\begin{explain}
\textbf{ARPIT} (Annual Refresher Programme in Teaching), launched by MHRD in November 2018 on the SWAYAM platform, provides online refresher courses for higher-education faculty. Completing it counts for career advancement, like a conventional refresher course.
\why{MOODLE is an open-source LMS; NPTEL provides engineering and science courses; e-PG Pathshala provides PG e-content.}
\end{explain}
\end{question}

\begin{question}{40}
{Which of the following is not a social network?\par}
\optsfour{Linkedin}{Instagram}{Twitter}{ebay}
\answer{D}{ebay}
\begin{explain}
LinkedIn, Instagram and Twitter (X) are social-networking platforms. \textbf{eBay} is an \emph{e-commerce} (online auction and shopping) website.
\end{explain}
\end{question}

\begin{question}{41}
{Which steps are to be taken for abatement of air pollution?\newline (i) Use a CNG Car.\newline (ii) Stop the use of fire crackers during festivals.\newline (iii) Plant more number of trees.\newline (iv) Not to burn agricultural waste.\newline Choose the correct options:\par}
\optstwo{(iii) only}{(i), (iii) and (iv) only}{(i), (ii), (iii) and (iv)}{(i), (ii) and (iv) only}
\answer{C}{(i), (ii), (iii) and (iv)}
\begin{explain}
All four help abate air pollution: CNG vehicles emit fewer pollutants (i); avoiding firecrackers prevents spikes in particulate matter (ii); trees absorb pollutants and dust (iii); and not burning crop residue (stubble) prevents severe smog (iv). The answer is (i), (ii), (iii) and (iv).
\end{explain}
\end{question}

\begin{question}{42}
{Ozone depletion in the stratosphere causes:\par}
\optsone{increase in greenhouse effect}{cooling of the earth atmosphere system}{increase of infrared radiation in the troposphere}{increase of ultraviolet radiation in the troposphere}
\answer{D}{increase of ultraviolet radiation in the troposphere}
\begin{explain}
The stratospheric ozone layer absorbs most of the Sun's harmful UV-B radiation. When it is depleted, \textbf{more ultraviolet radiation reaches the troposphere} and the Earth's surface, causing skin cancer, cataracts and crop damage.
\end{explain}
\end{question}

\begin{question}{43}
{According to report of NITI Aayog, which city from South India will have zero ground water by 2020?\par}
\optstwo{Kozhikode}{Thiruvananthapuram}{Chennai}{Mysuru}
\answer{C}{Chennai}
\begin{explain}
NITI Aayog's \emph{Composite Water Management Index} report (2018) warned that 21 Indian cities, including Delhi, Bengaluru, \textbf{Chennai} and Hyderabad, could run out of groundwater by 2020. Among the options, Chennai is the South Indian city named. It faced an acute ``Day Zero'' water crisis in 2019.
\end{explain}
\end{question}

\begin{question}{44}
{The Air Quality Index (AQI) of one of the city is 190. It falls under which category?\par}
\optsfour{Very poor}{Poor}{Moderate}{Good}
\answer{C}{Moderate}
\begin{explain}
India's AQI bands: Good (0--50), Satisfactory (51--100), \textbf{Moderate (101--200)}, Poor (201--300), Very poor (301--400), Severe (401--500). An AQI of 190 is \textbf{Moderate}.
\end{explain}
\end{question}

\begin{question}{45}
{Which among the following is an environmental friendly polymer?\par}
\optstwo{Poly styrene}{Poly ethylene}{Poly lactic acid}{Poly propylene}
\answer{C}{Poly lactic acid}
\begin{explain}
\textbf{Poly lactic acid (PLA)} is made from renewable resources such as corn starch or sugarcane and is \textbf{biodegradable} and compostable, so it is environment-friendly.
\why{Polystyrene, polyethylene and polypropylene are petroleum-based, non-biodegradable plastics.}
\end{explain}
\end{question}

\begin{question}{46}
{The purpose of the National Assessment and Accreditation Council (NAAC) in India is \blank{}.\par}
\optstwo{Granting research funds}{Approving Academic Programs}{Quality Assurance}{Establishing Universities}
\answer{C}{Quality Assurance}
\begin{explain}
\textbf{NAAC} (est.~1994 by UGC, Bengaluru) assesses and accredits higher-education institutions, so its purpose is \textbf{quality assurance}.
\why{Grants are given by UGC/HEGC, programmes are approved by statutory councils and universities, and universities are established by legislatures.}
\end{explain}
\end{question}

\begin{question}{47}
{Which of the following organizations is responsible for Technical Education in India?\par}
\optsfour{AIU}{AICTE}{NCTE}{AIIMS}
\answer{B}{AICTE}
\begin{explain}
The \textbf{AICTE} (All India Council for Technical Education) is the statutory body for technical education: planning, standards, approvals and so on.
\why{AIU is the Association of Indian Universities; NCTE regulates teacher education; AIIMS is a medical institute.}
\end{explain}
\end{question}

\begin{question}{48}
{`IMPRINT' is a program that addresses \blank{}.\par}
\optsone{Major Science \& Engineering Challenges in India}{Major Agriculture Challenges}{Improvement of Academic Research}{Major Language and Literacy Challenges}
\answer{A}{Major Science \& Engineering Challenges in India}
\begin{explain}
\textbf{IMPRINT} (IMPacting Research INnovation and Technology), launched in 2015 by MHRD and led by the IITs and IISc, is a national initiative to address \textbf{major science and engineering challenges} in ten technology domains (health, energy, water, environment, security and others).
\end{explain}
\end{question}

\begin{question}{49}
{Which institution secured first rank in over all category of NIRF ranking 2019?\par}
\optstwo{IISc, Bangalore}{IIT, Madras}{IIT, Delhi}{IIT, Mumbai}
\answer{B}{IIT, Madras}
\begin{explain}
In \textbf{NIRF 2019} the \textbf{IIT Madras} ranked 1st in the Overall category, followed by IISc Bengaluru and IIT Delhi.
\end{explain}
\end{question}

\begin{question}{50}
{10+2+3 pattern of Education system in India was recommended by \blank{}.\par}
\optstwo{Kothari Commission}{Radhakrishnan Commission}{Balvantrai Mehta Committee}{Janardhan Reddy Committee}
\answer{A}{Kothari Commission}
\begin{explain}
The \textbf{Kothari Commission} (Education Commission, 1964--66) recommended a uniform \textbf{10+2+3} pattern of education, which the National Policy on Education 1968 adopted.
\why{The Radhakrishnan Commission (1948--49) dealt with university education; the Balwantrai Mehta Committee (1957) dealt with Panchayati Raj; the Janardhan Reddy Committee (1992) reviewed the NPE 1986.}
\end{explain}
\end{question}
